
Semantic entropy (SE), a method published in \textit{Nature} by Farquhar et al. (2024) with over 1,000 citations, is widely used as a benchmark for uncertainty quantification (UQ) in language models. This paper reports that SE fails as an uncertainty estimator when applied to Llama-3-8B-Base on factual question answering: on 500 validation queries from TriviaQA (rc.nocontext), SE achieves an AUROC of 0.4735 (95\% CI: [0.4409, 0.5036]), compared to 0.6835 for token-probability (95\% CI: [0.6361, 0.7332]). On NaturalQuestions (open-domain validation split, 500 queries), SE achieves an AUROC of 0.5524 (95\% CI: [0.5121, 0.5977]) versus 0.6551 for token-probability (95\% CI: [0.5960, 0.7063]). The observed failure is consistent with a sampling degeneracy phenomenon: under N=10 stochastic sampling at temperature=1.0, 89.4\% of TriviaQA queries and 84.8\% of NaturalQuestions queries produce a single semantic cluster (K=1), collapsing SE entropy to near-zero for the majority of queries and eliminating its discriminative ability. Kernel Language Entropy (KLE) similarly fails, with AUROC of 0.2642 on TriviaQA and 0.3753 on NaturalQuestions. SelfCheckGPT-BERTScore yields a constant AUROC of 0.5000 on both datasets. SelfCheckGPT-NLI achieves AUROC of 0.6862 on TriviaQA but only 0.4508 on NaturalQuestions, indicating dataset-dependent performance. The paper introduces \texttt{degenerate_fraction}—the proportion of queries for which all N samples cluster into a single semantic equivalence class—as a diagnostic for SE validity, and provides a complete benchmarking infrastructure of approximately 2,222 lines of Python code with checkpointed generation and bootstrap confidence intervals. The 70B-scale experiment was initiated but did not complete before the time of reporting; the 8B gate results alone are sufficient to falsify the existence claim that SE exceeds token-probability for base models on these benchmarks.

